% Section 4 (Methodology) continued — Simulation Model through Experimental Design
% Sources: src/simulation.py, src/slotting.py, src/config.py, src/kpi.py,
%          src/validate.py, src/analyze.py, ASSUMPTIONS.md §15/21/23/24,
%          report/sections/gen/facts.json

\subsection{Simulation Model}\label{sec:simmodel}

\subsubsection{Modelling Approach}

We built the simulation using Python's SimPy~\cite{simpy} library, following
the process-interaction world-view described in Banks et~al.~\cite{banks2010}
and Law~\cite{law2015}. In process-interaction discrete-event simulation each
entity moves through a sequence of activities and queues; the simulation clock
advances only to the next scheduled event, so idle time costs nothing
computationally.

The primary entity is a \emph{kit order}: a release of $n$ distinct component
materials that must be picked from the warehouse and assembled at a production
line's kitting point. Each order enters a queue for an \emph{operator}
resource (capacity~8), and the operator then walks to each pick location in
turn. When a pick location is at level~3 or above --- level~3 is the first shelf
not reachable by an operator on foot from the kit corridor (\texttt{FAST\_MOVER\_MAX\_LEVEL\,=\,3},
the three lowest levels are kit-accessible) --- the operator must also request a
\emph{reach truck} (RT, capacity~7). Picks from the automated Kardex carousel
system are handled by a separate Kardex resource pool (capacity~4). A
per-position SimPy resource lock (capacity~1) prevents two operators from
attempting the same pallet slot simultaneously, which would be physically
impossible.

After delivering the last item to the kitting point, the RT returns to its
depot as a detached SimPy process that keeps the truck resource busy until it
arrives --- so the operator is released at drop-off but the truck is not
available for the next request until the return leg completes. This detail
matters: without it the model understates RT utilization.

\begin{figure}[htbp]
  \centering
  \includegraphics[width=0.82\textwidth]{sim_flowchart_simple.png}
  \caption{Overview of the life-cycle of one kit order in the simulation.
           The full expanded flowchart with branching conditions is given
           in Appendix~\ref{fig:flow_detailed}.}
  \label{fig:simflow}
\end{figure}

Figure~\ref{fig:simflow} shows the order life-cycle at a summary level.
The model is \emph{terminating}~\cite{law2015}: each run covers five
eight-hour working days, matching the ZWM92 dispatch profile. A 30-minute
warmup and 30-minute cooldown window trim edge transients from the KPI
window, so no steady-state warm-up analysis is needed.

\subsubsection{Order Generation}

The simulation uses a \emph{distribution-driven} input model in the
Arena~\cite{law2015} tradition: distributions are fitted once to four
months of ZWM92 warehouse management dispatch data (167,784 rows, 40,804
kit-orders across nine product families, January through May~2026) and each
replication draws independent samples. The raw trace is not replayed.

When we analyzed the ZWM92 log, we found that kit-orders are not independent
arrivals. Among the 18,164 timestamped orders, 3,885 share an identical
timestamp, forming batches with a mean of about 4.68~kits (median~2,
maximum~79). Each batch corresponds to a single production order releasing
all its kits at once. Treating the Exp(5.36~min) fitted inter-batch gap as
a per-order inter-arrival time would have under-loaded the model by roughly
4.5$\times$, producing about 87~simulated orders per day instead of the
real average of about 396~orders per active working day (40,804 orders over
103 active calendar days). We therefore switched to a batch arrival model:
one arrival event generates one batch, not one kit.

Inter-batch gaps are drawn from an empirical pool of 3,604 within-shift
positive gaps; the coefficient of variation of 2.04 rules out an exponential
fit, so we kept the empirical pool. Batch size is drawn from the 3,885
observed run lengths. The production line is drawn once per batch from a
categorical distribution over the nine ZWM92 family weights --- 99.9\% of
multi-kit batches in the log are single-line, which makes this a safe
simplification. Each kit in the batch draws its item count from the empirical
distribution of distinct components per kit (mean~4.0, clipped at~50), and
each item is drawn from the material-weight distribution for that line.

Calibration: the batch driver's theoretical expected output is approximately
419~orders per day ($480\,\text{min} \div 5.36\,\text{min/batch} \times 4.68\,\text{kits/batch}$),
against the ZWM92 target of about 396 orders per active day --- a $+5.7\%$
difference, inside the $\pm$10\% acceptance band. The validation module
checks this automatically on every run and writes a
\texttt{daily\_volume\_check} result to the output.

\subsubsection{Travel, Picking and Service Times}

\paragraph{Routing.}
All travel in the simulation is \emph{corridor-aware rectilinear}: rack
bodies are treated as blockers, candidate paths snap to a grid of safe lines
derived from rack edges plus a 60-cm clearance margin, and the shortest
obstruction-free path wins. Because a Manhattan L-bend is a proven lower
bound, any clear L-path is returned immediately without further search. Path
results are cached at the class level (keyed to a geometry fingerprint), so
the routing cost is paid only once per layout~\cite{bartholdi2014}. Operators
walk at 50~m/min; RTs travel at 100~m/min.

\paragraph{Operator picks.}
When a material's slot is in the three lowest levels and the rack has a kit
corridor, the operator picks directly from the kit-corridor face. At level~3 or above
the operator must wait for an RT. A fallback manual pick is allowed at
level~2 or below from the RT-aisle face (no truck needed), but carries a time
penalty modelled as the corridor-traversal mean from the F400 video study
(0.102~min per pick --- a conservative upper-bound proxy that overstates the
true marginal penalty by roughly 30\%; see ASSUMPTIONS~\S19). Picks are
assigned in a greedy nearest-first sequence over routed distances; for
materials that occupy multiple SAP bins, the simulation picks from the
nearest available slot.

\paragraph{Reach-truck picks.}
The RT travels from its depot (south-west of the rack block) to the pick
face, lifts at 0.25~min per level, and executes the pick/place (mean
0.11~min) before carrying the pallet to the kitting point. The operator is
released at drop-off; the truck then returns to the depot under its own
process, holding the RT resource until it arrives back.

\paragraph{Kardex picks.}
All materials classified as Kardex items in the SAP master (2,872~materials
in the model) bypass the pallet rack entirely. The simulation requests the
Kardex resource, pays one carousel rotation (0.4~min, book value), then
executes per-item pick time (0.113~min per item, mirroring the operator
pick constant) inside a single resource hold. The four Kardex units share one
SimPy resource pool with capacity~4.

\paragraph{Stochastic variability.}
All pick and travel times are stochastic when \texttt{STOCHASTIC\_PICK\_TIMES}
is enabled. Each duration is drawn from a lognormal distribution whose mean
matches the F400-measured constant and whose log-scale standard deviation
$\sigma = \sqrt{\ln(1 + \mathrm{CV}^2)}$ is moment-matched from the observed
coefficient of variation in the video study: operator pick $\sigma = 1.245$
(pooled rf\_scan + manual\_pick, $\mathrm{CV} = 1.93$, $n = 900$); RT
pick/place $\sigma = 1.047$ ($\mathrm{CV} = 1.41$, $n = 69$); manual-pick
penalty $\sigma = 1.279$ ($\mathrm{CV} = 2.03$, $n = 157$); Kardex pick
$\sigma = 1.245$ (mirrors operator; no Kardex video available).

\subsubsection{Model Assumptions}\label{sec:assumptions_condensed}

The key assumptions are listed below; a complete register appears in
Appendix~\ref{sec:app_assumptions}.

\begin{enumerate}
  \item \textbf{Single position per material.} Each material is placed at
        one pallet slot (or at all its decoded SAP bins when multi-bin). The
        model does not simulate replenishment: pallets never deplete and no
        replenishment picks occur. Because this holds for all six policies,
        the comparison is internally fair.
  \item \textbf{F400 timing extrapolated.} The video time-motion study
        covers only the F400 production line, which accounts for about 48,000
        of the 167,784 ZWM92 dispatch rows. F400-measured constants are
        applied to all nine families. The sensitivity analysis bounds the
        risk of this assumption.
  \item \textbf{No aisle congestion between agents.} Operators and reach
        trucks do not block each other in shared corridors; each queues for
        the resource and then travels at full speed. Congestion effects are
        listed as a model limitation.
  \item \textbf{Single-shift operation.} ZWM92 shows orders in a
        07:00--16:00 window. The model runs as a continuous flow (no
        overnight gap, no within-day breaks) with warmup and cooldown trim.
        Multi-shift dynamics are out of scope.
  \item \textbf{Distribution-driven, not trace-replay.} Each replication
        draws independent samples from fitted distributions; the raw ZWM92
        sequence is not replayed. This is the standard Arena-style
        approach~\cite{law2015}.
  \item \textbf{Fit and validation share the same log.} Distributions were
        fitted to, and the chi-square validation expected vectors were drawn
        from, the same ZWM92 dataset. These tests are therefore
        internal-consistency checks rather than holdout tests. External
        corroboration comes from the F400 video timing, the daily-volume
        calibration, and face-validity checks against site photos and the
        AutoCAD floor plan.
\end{enumerate}


\subsection{Storage Assignment Policies}\label{sec:policies}

Six storage assignment policies are compared. All six place the same 5,941
active materials in the same physical warehouse and are exposed to the same
demand process. The only difference between experiments is where each
material sits.

Our initial proposal focused on classifying materials by demand recency,
which naturally led us to the frequency-ranked approaches first --- the
usage-based and travel-distance policies below. As we developed the model
and identified the congestion mechanism described in Section~\ref{sec:vv},
the proposed policy evolved into the line-aware variant.

\paragraph{Policy~1: Current SAP placement (Baseline -- Actual SAP).}
The ground-truth baseline. Storage bin codes are read directly from the SAP
\texttt{özet} material master (\texttt{BRA-02-02} $\to$ rack~A, bay~2,
position~2). Materials whose bins decode to a free pallet slot are placed
there; Kardex-flagged materials are routed to the Kardex resource; the
remaining materials without a decodable SAP rack bin fall back to an
FMR-heuristic pool. Of 5,941 active materials: 750 are placed at their
exact SAP rack bin, 2,872 are Kardex-routed (policy-invariant across all
policies), and the rest fall back to the heuristic pool. This baseline is
the company-relevant comparator and the primary contrast partner for
hypothesis testing.

\paragraph{Policy~2: Capacity heuristic (Baseline -- Heuristic).}
Materials are sorted by the existing SAP FMR class (F~= fast, M~= medium,
R/D~= rare/dead) and filled into position pools in that order: fast movers
to the lower-level pool, medium movers to mid-level, and rare/dead movers to
upper levels. No actual demand data are used. This policy fills capacity
without considering per-material frequency; any data-driven policy should
outperform it, and it is the heuristic-pool fallback for Policies~1 and~6.

\paragraph{Policy~3: Usage-based ABC.}
Materials are reclassified using ZWM92 pick counts rather than the SAP FMR
flag: class~A (top 80\% of cumulative picks) goes to the fast-mover
positions, class~B (80--95\%) to mid-level, and class~C (bottom 5\%) to
upper levels. The thresholds follow the Pareto principle
applied to real warehouse demand~\cite{frazelle2002, dekoster2007}.

\paragraph{Policy~4: Double ABC.}
This cross-tabulates the SAP FMR class (frequency axis) against the
usage-based ABC class (volume axis), yielding twelve combined classes (AF,
BF, AM, \ldots, CR). Materials are sorted by a fixed priority list --- AF and
BF take the best positions first, followed by AM, CF, BM, AR, CM, BR, and
CR --- and filled into the level pools in that order~\cite{flores1986}. The
intent is to balance volume and frequency jointly.

\paragraph{Policy~5: Travel-distance Optimized.}
Materials are ranked by their ZWM92 pick count (falling back to the SAP
consumption column when the ZWM92 cache is unavailable) and then assigned
greedily: the highest-frequency material gets the position closest (by
routed distance) to the central kitting centroid, the next highest gets the
next closest, and so on. This is a standard nearest-to-dock heuristic in the
warehousing literature~\cite{hausman1976, bartholdi2014}. As we show in
Section~\ref{sec:results}, it produces the worst lead times of all six
policies under realistic load, because it routes all high-frequency
materials to positions near the central kitting point regardless of level ---
and those positions are predominantly RT-served, so every pick queues a reach
truck.

\paragraph{Policy~6: Line-aware Slotting (proposed).}
This policy directly addresses the RT-congestion problem exposed by Policy~5.
Materials are sorted by ZWM92 pick frequency and assigned in three passes.

First, each material's \emph{origin} is set to the kitting point of its
dominant production line, read from the \texttt{production\_lines} array in
the warehouse layout. Four lines --- F400, SM6-36, PREMSET, MCSET --- have
confirmed kitting points covering about 76\% of ZWM92 pick volume; the
remaining lines fall back to the central centroid. Every warehouse position
is pre-scored by its routed distance from that origin. Among free positions,
the policy selects those that do \emph{not} require a reach truck: slots in
the lower three levels of a rack with a kit corridor, or at a
manually-reachable level in an RT aisle. Fast-moving materials therefore land
on hand-reachable levels, largely bypassing the RT queue. Because each
line's demand is drawn to its own corridor, load spreads across multiple
corridors rather than concentrating on one.

Second, if the no-RT space for a line is exhausted, the policy falls back to
the nearest RT-served position from the same line origin.

Third, zero-pick materials (those absent from ZWM92 entirely) are placed last,
in upper-level reserve positions, so prime slots stay available for active
SKUs.

Why this outperforms the other policies is discussed in
Section~\ref{sec:results}.

% --------------------------------------------------------------------------
\subsection{Performance Measures}\label{sec:kpis}

We report seven performance measures. The four primary KPIs --- kitting lead
time, total waiting time, reach-truck utilization, and daily throughput ---
were set together with our supervisor; walking distance and the remaining
measures support the causal explanation.

\emph{Kitting lead time} (minutes per order) runs from the moment the kit
order joins the operator queue to the moment the last item is delivered to the
kitting point. It is the end-to-end measure the warehouse manager cares about
most: a longer lead time means the production line waits longer for its parts.

\emph{Waiting time} (minutes per order) is the combined time an order spends
in the operator queue and the RT queue. It isolates resource contention from
the picking activity itself, so it helps explain why lead times differ across
policies.

\emph{Reach-truck utilization} is the fraction of cumulative RT-resource
capacity consumed over the full simulation window. Very high utilization
signals a congestion bottleneck; it is the causal variable we use to explain
Policy~5's poor lead time.

\emph{Throughput} (kit-orders per simulated day) measures overall output rate.
Under the experimental conditions here, all policies achieved similar
throughput, which means lead-time differences are driven by waiting, not by
capacity shortage.

\emph{Walking distance per order} (metres) is the total rectilinear path
length the operator walks during one kit order. As the results show, reducing
walk distance in isolation is not sufficient to improve lead time when doing so
shifts picks to RT-served positions.

\emph{Preparation time} (minutes per order) is the operator-active duration
after the resource is granted: walk plus pick time, excluding queue wait. It
complements lead time by showing how much of the delay is active effort versus
idle waiting.

\emph{Operator utilization} is defined analogously to RT utilization:
cumulative operator-busy time divided by total operator capacity over the
simulation window.


% --------------------------------------------------------------------------
\subsection{Verification and Validation Approach}\label{sec:vv}

We follow Sargent's framework~\cite{sargent2013}, which distinguishes
\emph{verification} (does the program correctly implement the conceptual
model?) from \emph{validation} (is the model an accurate enough
representation of the real system for our purpose?).

\paragraph{Verification.}
Code-level tests cover the SAP bin decoder, Kardex routing logic, and
resource-accounting identities. A ten-question sanity gate
checks that total pallet capacity equals
the PDF-stamp value of 3,203; every placed material has a valid position;
pick counts sum correctly; the CRN seed pair produces identical arrival
sequences across policies at the same replication index; and several other
internal consistency conditions. The 3,203-position canary must pass before
any run is reported. Deterministic seed replays confirm bit-identical output
across Python sessions.

\paragraph{Face validation.}
Layout and material flow were validated against the physical facility through
a 3D interactive viewer we developed (available via the project web page),
walkthroughs with site photographs, and two supervisor meetings. Floor
dimensions, rack positions, corridor widths, and the kit-corridor/RT-aisle
alternation pattern were confirmed against the AutoCAD floor plan and eleven
rack-layout PDFs. Kit-corridor widths of 2~m and RT aisle widths of
300--320~cm were measured directly from the CAD dimension annotations.

\paragraph{Operational validation.}
We compare the SAP-baseline simulation's per-rack pick distribution against
ZWM92 actual pick counts using a chi-square goodness-of-fit test, and run a
per-material paired $t$-test comparing simulated daily pick rates against
observed ZWM92 rates for the 925~materials that could be paired. Both tests
are run by an automated validation script; results are summarized in
Table~\ref{tab:validation} and discussed in Section~\ref{sec:results}.

Both tests reject the null hypothesis, as we expected. The chi-square asks
whether the simulation reproduces SAP's \emph{exact} bin assignments; it
cannot, because each policy rebuilds placements from scratch and no policy
can reconstruct the bin decisions that SAP accumulated over years of manual
and semi-automated putaway. The relevant measure is therefore Cram\'{e}r's
$V$, which quantifies effect size rather than significance~\cite{sargent2013}.
The headline $V = 0.22$ (restricted test: both the simulated picks and the
expected vector are scoped to the 750 materials with decoded SAP rack bins,
giving an apples-to-apples comparison). The per-material $t$-test also
rejects, but the mean log-ratio of 0.32 --- simulated daily rates about 37\%
above ZWM92 actuals --- is consistent with the batch model slightly
overshooting its calibration target. Both the scope mismatch and the
overshoot are disclosed as limitations.

A daily-volume check verifies that the batch driver stays within
$\pm$10\% of the ZWM92 target (about 396~orders per active day); the
single-run validation produced about 378~orders per day, a $-4.6\%$ error,
inside the acceptance band. A replication CI check confirms that the 95\%
confidence interval on throughput across 20~replications is consistent with
the target volume. Both checks are logged automatically on every run.

Because the ZWM92 log was used both to fit the input distributions and to
build the validation expected vectors, the chi-square and $t$-tests are
internal-consistency checks, not holdout tests. External corroboration comes
from the F400 video timing, the daily-volume calibration, and the face-validity
checks described above. All numerical outcomes are reported in
Section~\ref{sec:results} (Table~\ref{tab:validation}).


% --------------------------------------------------------------------------
\subsection{Experimental Design and Statistical Analysis}\label{sec:doe}

\paragraph{Replications.}
Each policy was run for $N = 20$ independent replications, each covering five
simulated working days (480~minutes per day). Twenty replications exceeds the
advisor's minimum criterion and gives enough degrees of freedom for the ANOVA
and pairwise $t$-tests.

\paragraph{Common random numbers.}
Across policies, replication $i$ uses seed $42 + i$ for the arrival stream
and seed $(42 + i) + 1{,}000{,}003$ for the service-time stream. Because the
seed depends only on the replication index, every policy sees exactly the
same arrival trajectory at the same replication, implementing common random
numbers (CRN)~\cite{law2015}. CRN reduces the variance of policy-difference
estimators by cancelling demand-side noise; all six policies are exposed to
the same sequence of kit orders within each replication.

The two-stream design (arrival RNG and service RNG separated by a large
offset) is important: a single shared stream would let a policy that makes
one extra service-time draw shift every subsequent inter-arrival time, silently
breaking CRN. With separate streams, the demand trajectory is identical
across policies regardless of how many service draws they make.

\paragraph{KPI export.}
Each replication appends one row to a tidy results table (policy,
replication index, and every KPI as columns), giving $6 \times 20 = 120$
rows in the format Minitab expects, which we used for cross-checking.

\paragraph{Hypothesis tests.}
Three layers of inference are applied.

\emph{One-way ANOVA} tests whether all six policy means are equal for each
KPI, with rejection at the $\alpha = 0.05$ significance level. For lead time the $F$-statistic is $F(5, 114) = 47.07$ with
$p < 10^{-25}$, so the null is clearly rejected.

\emph{Tukey HSD post-hoc tests} provide 95\% family-wise confidence intervals
and adjusted $p$-values for all $\binom{6}{2} = 15$ policy pairs, identifying
which specific differences are significant while controlling the family-wise
error rate.

\emph{CRN-paired $t$-tests} compare each alternative policy against the SAP
baseline by pairing replication~$i$ of the alternative against replication~$i$
of SAP. Pairing on replication index removes demand-side variance and gives
tighter confidence intervals than an unpaired test~\cite{law2015}. The 95\%
CI half-width uses the $t(19)$ critical value. Mean differences, 95\% CIs,
and $p$-values are reported in Section~\ref{sec:results}
(Table~\ref{tab:paired_t}).

All numerical results and the full ANOVA, Tukey, and paired-$t$ tables are
presented in Section~\ref{sec:results}.
